% !PRINTED 189
\weberpage{189}

Es sei schliesslich noch bemerkt, dass, anstatt der \gesperrt{Stellungen} des
Körpers, auch die Drehungen selbst als Elemente der Gruppe
aufgefasst werden können.

\begin{center}
\S.\,49.

\gesperrt{Lineare gebrochene Substitutionen.}
\end{center}

Die Gruppe der orthogonalen ternären Substitutionen ist, wie
jetzt nachgewiesen werden soll, isomorph mit der Gruppe der
linearen gebrochenen Substitutionen einer Veränderlichen, oder der
binären Substitutionen der Verhältnisse (\S.\,38).

Wir bezeichnen
eine ternäre orthogonale Substitution mit
\weq{1}{(y_1,\ y_2,\ y_3) = A\,(x_1,\ x_2,\ x_3),}
\weq{2}{y_1^2 + y_2^2 + y_3^2 = x_1^2 + x_2^2 + x_3^2.}
Hierauf wenden wir zunächst nach \S.\,37, (21) die
Transformation durch die feste (nicht orthogonale) Substitution
\weq{3}{L = \left\{\begin{array}{ccc}
i, & 0, & 0 \\[4pt]
0, & \dfrac{1}{\sqrt{2}}, & \dfrac{1}{\sqrt{2}} \\[6pt]
0, & \dfrac{i}{\sqrt{2}}, & \dfrac{-i}{\sqrt{2}}
\end{array}\right\},\quad
L^{-1} = \left\{\begin{array}{ccc}
-i, & 0, & 0 \\[4pt]
0, & \dfrac{1}{\sqrt{2}}, & \dfrac{-i}{\sqrt{2}} \\[6pt]
0, & \dfrac{1}{\sqrt{2}}, & \dfrac{i}{\sqrt{2}}
\end{array}\right\}}
an, worin $i = \sqrt{-1}$ ist, deren Determinante den Werth 1 hat,
und setzen
\weq{4}{\begin{aligned}
(y_1,\ y_2,\ y_3) &= L\,(y_1',\ y_2',\ y_3')\\
(x_1,\ x_2,\ x_3) &= L\,(x_1',\ x_2',\ x_3')' % [sic] source prints a stray prime after the closing parenthesis
\end{aligned}}
\weq{5}{(y_1',\ y_2',\ y_3') = A'\,(x_1',\ x_2',\ x_3'),}
worin
\weq{6}{A' = L^{-1} A L}
gleichfalls eine lineare Substitution bedeutet, deren
Determinante $|A'| = |A|$, also gleich $\pm 1$ ist. Die Relation
(2) geht durch die Substitutionen (4) in
\weq{7}{-y_1'^2 + 2\,y_2'\,y_3' = -x_1'^2 + 2\,x_2'\,x_3'}
über, woraus sich sechs Relationen zwischen den Co\"efficienten
von $A'$ ergeben, die wir aber hier nicht aufstellen wollen, da
wir die Substitution $A'$ auf andere Weise einfacher finden
können.
